Full soft targets yielded lower selective disagreement than the ambiguity-preserving uniform-minority control under this matched protocol. This is a test of whether minority allocation supplies added selective value beyond dominant-confidence-preserving softening, not a claim that soft-label learning itself is new. The tie-aware condition directly probes a second alternative explanation rather than adding model capacity. The selected and final-epoch results delimit the dependence on the stopping rule; neither warrants universal superiority.

The size of the control contrast matters: Uniform lowers selected-checkpoint disagreement relative to Hard by 2.23 pp, compared with 2.61 pp for Soft. Thus, the simpler target reproduces much of the observed hard-to-soft difference; detailed minority allocation adds a modest residual advantage under this protocol. This is an endpoint comparison, not a percentage of improvement causally explained by softening.

The vocabulary diagnostics explain why a high supported-class confidence cannot by definition encode all complete-vote risk. Unsupported mass and supported ambiguity are separately observable in the held-out annotations, while MSP omits the former even at the ideal supported-distribution optimum. Nevertheless, observed risk gaps also contain excess category-choice error. The decomposition is useful precisely because it prevents attributing a total change to a single mechanism without evidence.

The conservative grid comparator adds a specified finite-sample margin, at a potential coverage cost. Its assumptions have not been established for deployment or person-level independence in this benchmark, so empirical compliance is not proof of a general guarantee. Calibration half-splits are similarly diagnostic rather than a new validation cohort. The practical recommendation is to report achieved risk, coverage, annotation vocabulary, and accepted-set composition together, instead of treating a nominal confidence policy or majority-only F1 as sufficient reliability evidence.

